% ============================================================
%  Poul Martin Møller, En Dansk Students Eventyr (1824)
%  TRANSLATION (English)
%  Translator: [HH]
%  Selection: The Licentiate episode
% ============================================================

\documentclass[12pt]{article}
\usepackage[utf8]{inputenc}
\usepackage[T1]{fontenc}
\usepackage[english]{babel}
\usepackage{libertinus}
\usepackage{libertinust1math}
\usepackage[protrusion=true,expansion=true]{microtype}
\usepackage[margin=1.4in]{geometry}
\usepackage{setspace}
\usepackage{fancyhdr}
\usepackage[colorlinks=true,linkcolor=black,urlcolor=black,
  pdftitle={A Danish Student's Story (excerpt)},
  pdfauthor={Poul Martin Møller}]{hyperref}
\setlength{\headheight}{15pt}
\pagestyle{fancy}
\fancyhf{}
\fancyhead[LE,RO]{\thepage}
\fancyhead[RE]{\textit{A Danish Student's Story}}
\fancyhead[LO]{\textit{The Licentiate}}
\setlength{\parindent}{1.5em}
\setlength{\parskip}{0pt}
\onehalfspacing

\title{\textbf{A Danish Student's Story}\\[0.5em]
  \large Excerpt: The Licentiate}
\author{Poul Martin Møller\\[0.4em]
  \small Translated by [HH]}
\date{First published 1824; revised text in
  \textit{Efterladte Skrifter}, vol.~3,\\
  ed.\ F.C.\ Olsen. Copenhagen: C.A.\ Reitzel, 1843}

\begin{document}
\maketitle

\section*{The Licentiate}

\noindent
Three years ago, the family at Ravnshøi and a number of their
well-to-do neighbors, who have no publicly appointed physician in
their vicinity, had joined together to engage a doctor at their own
expense. I resolved to seize this opportunity to bring my dear
kinsman into some kind of occupation. For he is himself far too lame
of will to carry out any decision. I put a letter in his pocket, set
him on a goods wagon, and thought the matter perfectly in order. But
when a year or so later I made my way out there to see how he was
managing in his new position, and for this purpose had myself taken
to the local physician, I found to my great astonishment that another
doctor was practicing his art in that part of the country. I now
resolved to go back to the city to seek out the Licentiate there if
possible. On the way I spent the night at the Mill Inn, and there,
while I sat in the evening dusk with the jolly landlord watching the
fire in the tile stove, I saw a tall, crooked figure steal along the
wall and disappear through one of the doors. This figure I shall
describe; for in so doing I describe my cousin at the same time. He
stands three and a half ells tall, and out of shame at this
supernatural height, which is made all the more conspicuous by his
slenderness, he bends his back so far that his head falls below the
sight-line of an ordinary human being. Nature has eased him in this
posture by furnishing his spine with two joints, so that it forms,
especially when he is seated, three lines making two obtuse angles
with one another. Otherwise there is nothing remarkable about his
exterior, save for a clubfoot, which makes it necessary for him to
wear a lace-up boot on his right foot.

--- What sort of man was that? I asked the landlord.

--- He is one of my guests, he replied; he arrived at my house a year
ago and appeared to be in great haste, for he ordered hot water for
the following morning at four o'clock, so as to continue his journey
in good time. But the next day he cancelled the carriage. He repeated
this every day for the first half year; but now it seems he has
settled in here at my house. He is a good fellow who never gets in
anyone's way; but I believe, to tell the truth, he is not quite right
in the head. His learning has driven him out of his senses.

--- Choose your words carefully, I said, for he is my own first cousin
on my mother's side.

--- Ah, he said, then you must know best yourself how he is. He is
terribly unsociable. Yesterday he had a ladder set to his room window,
so that he could come down into the garden and take his customary
exercise without going through the taproom. Sometimes he lies in a
stupor the whole day; sometimes he paces up and down the floor and
begins to hold forth to himself. But it is quite superfluous for me
to tell you such things. You must surely know his peculiarities best
yourself. He gives us no trouble, as I said. We have grown so
accustomed to seeing him lumber about the house that I even believe we
should miss him if he went away. My daughter Marie takes great care
to see that he gets his food and drink in good time, for otherwise I
believe he would forget himself and simply starve to death. I tell
her often enough: let us see for once whether nature itself cannot
remind him. He is so tall, after all, that he ought to be able to
make some demand upon himself. But she is too soft-hearted to attempt
the experiment. ---

I sat up until well into the night writing letters, and all the while
heard him continually stirring in his room. Now he would play a few
minutes on his flute, now he would pace about his floor in silence
with heavy steps. I went out into the garden to peep through his
windows. The candle on his table had burned far down into the socket,
for in his unearthly thoughts he had forgotten to push it up. The
flame rose and fell in the tin tube, which was quite blue from the
heat, and the melted tallow spread its stench through the whole room.
Half lying, half sitting, he was ensconced in a very low armchair,
with his long legs stretched out on the floor. He was sunk in
contemplation and stared fixedly before him through his spectacles.
Both his hands, which still held the flute, lay entirely motionless
in his lap. I watched him for several minutes without his making any
sign of moving a limb. Had I not seen the burning gaze through the
spectacles, I would have taken it for a successful wax cast of my
cousin the Licentiate. At last I called to him through the open
window: Good evening, cousin! --- He started, without being brought
out of his petrified posture. I then crawled through the window with
the help of the recently installed ladder, and by the force of my
fists brought him out of his stupor. He rose with a good many
embarrassed bows and stammered some incomplete compliments, by which
he expressed his thanks that I had so willingly taken the trouble of
helping a little with his plan of life, and made a number of
apologies for not having followed the course I had prescribed for
him. His thoughts move so slowly that he sometimes lacks the presence
of mind to answer the simplest question. He is especially
incomprehensible when he has lately been rapt in his metaphysical
reveries. At such times there is such a divorce between his words and
his meaning that he is unable to produce any complete sentence; all
his thoughts are stillborn. When some equilibrium had been restored
to his scattered mind, he made his confession, which~--- with the
omission of his frequent pauses and failed periods~--- ran as follows:

--- The day after I arrived here, it was my firm intention to go to
Ravnshøi; but I found it impossible to locate my pocket-scissors, and
you would hardly require me to present myself to a household with
unclipped fingernails. When the following day I grasped a firebrand
from the chimney to light my pipe, I burnt a large blister on my
right forefinger. With such a physical pain, you will readily
understand that I lacked the free circulation of mind required when
one is to introduce oneself with a suitable greeting to a real
Councillor's wife. When on the third day the miller's wife, following
my instructions, was dressing my injured hand, and wished to burn the
thread over with a lit taper, seven large drops of tallow fell on my
black coat. When the first drop fell I foresaw that more would follow;
but I was so overcome with gratitude for her solicitude that I had
not the courage to draw her attention to it. A long time then passed
before I managed to procure some turpentine to remove these spots.
Meanwhile my finest collars had become soiled, and several weeks
elapsed before I had them sent to a washerwoman. With such small
misfortunes, sent down upon me by a malicious demon, a year passed
without it being possible for me to get to Ravnshøi, which lies only
half a mile from this inn.

--- But good Lord, I exclaimed, how is it conceivable that on not a
single one of those 365 days you could muster the heroism to make
that herculean effort in spite of all obstacles?

--- That, he replied, is what you and everyone who has the good
fortune to be able to act without reflecting on their life are always
saying. Had I been able to discern any sufficient reason why the
journey should take place on one day rather than another, it would
have come about. But the misfortune consists precisely in this, that
I see far too clearly the truth that any action can without serious
harm be postponed a day.

--- But there you are simply wrong, cousin, I replied; nothing in the
world that can be done today ought to be put off till tomorrow.

--- If you deny me that, he took up again, then you must necessarily
grant me that it can possibly just as well be done a minute or a
second later; and this accursed possibility persists as long as one
draws breath. No one can prove to me that a postponement lasting a
second or a third\footnote{In the old sexagesimal division of time, a
  \emph{secund} is what we now call a second (one sixtieth of a
  minute), and a \emph{terts} is one sixtieth of a second.} should
have ruinous consequences, and because no one can prove it to me, my
life comes to a complete standstill. My infinite brooding over it
means that I accomplish nothing. Furthermore I come to think about my
thoughts on the matter~--- indeed, I think about the fact that I
think about it, and divide myself into an infinitely regressing series
of selves that observe one another. I do not know which self is to be
regarded as the genuine one, and the moment I stop at one, it is
again a self that stops at it. I grow dizzy and overcome with
vertigo, as though I were staring down into a bottomless abyss, and
the thinking ends with my feeling a dreadful headache.

--- Dear cousin, I said calmly, with such broodings you will torment
yourself as fruitlessly as if you turned your Goliath frame inside
out and tried to walk on your head with your long legs in the air. I
cannot help you at all with sorting out your many selves. It lies
entirely outside my sphere of competence, and I should have to be or
become just as mad as you if I were to engage in your superhuman
reveries. My way is to keep to tangible things and walk along the
broad highway of common sense; that is why my selves never get into a
muddle. Besides, I see well enough that the fault is mine. I ought
to have led you straight to your destination, opened the Councillor's
wife's door, and pushed you into the room in such a fashion that you
would have found yourself obliged to seize upon whatever self came
first to hand, in order to bid the lady and her good family good day.
But tell me now, what have you been doing here for a full year?

--- I have, he said, been collecting materials for my great work on
the physical and intellectual nature of man.

--- Oh, I replied, are you still thinking of that? Ten years ago you
had the same fixed idea. It must be becoming a thoroughly substantial
piece of work. I thought you were studying mineralogy, since I saw
the great quantity of stones here on your floor.

--- They are here solely for the sake of the work, he said. I have
observed that thoughts flow best from the pen when one possesses
entirely suitable writing implements. To have good pens, one must
have good knives, and to sharpen them, one must possess really choice
whetstones. Therefore, before proceeding further, I ordered various
kinds of stone, and have indeed, as you see, had to penetrate rather
deeply into mineralogy as a preparatory study.

--- To which I replied: This seems a thoroughly rigorous method,
provided you add to it a diligent study of ornithology in order to
solve the problem of which birds' feathers yield the best quill-cases,
and also sacrifice a few years to pedagogy in order to discover how
a boy is best educated so as to become in time a capable and upright
knife-grinder.

--- Strictly speaking, that is not taking it too far, he said. You
mock because you do not understand me. These experiments in
facilitating intellectual activity through the convenience of material
means also furnish contributions to the work itself, just as the
observation of how I came to undertake the investigation constitutes
a new chapter. Furthermore~---

--- Now stop there, cousin, I interrupted him, otherwise our thoughts
will run into a tangle again, as when we lost our way among our
selves. One should not row so far out that one loses sight of land.

In order to extricate myself and my cousin from the tangled cobweb of
argument into which he is apt to draw us until I with my Jutland
understanding am at last at a complete loss, I set about examining
the state of his clothes; for in every practical matter he is utterly
helpless. I found that his wardrobe, with the assistance of the
benevolent Marie, had been put in better order than I had expected.
While I was occupied with this, I noticed that by keeping his coat
buttoned he was giving his person a broader cross-section, so as to
block my view of a certain corner. Whenever I tried to cast my eye
in that direction he would trip around in front of my line of sight
like a goose which with outspread wings seeks to defend its newly
hatched brood against the dairymaid's hands; but she kicks the
short-sighted goose aside with her wooden shoes and counts all the
goslings in her flock in order to bring them in for the night.
In like manner I, having detected the mischief, seized the desperate
Licentiatus by the arm and said: Stand aside, cousin, I must see
what contraband you have hidden here in the corner. Here I discovered
a pile of mouldy collars and neckcloths which he had very carefully
managed to conceal from Marie's watchful eye. With a ruler I stirred
up the ill-kept linen and in so doing disturbed an entire tribe of
earwigs in their domestic felicity. The fleeing swarm of insects
escaped for the most part from my murderous pursuit by means of their
numerous legs and supple bodies; but I cried out in indignation:

--- Dearest cousin! Why in all the world do you conceal your failings
from me? That you hide your slovenliness from the pretty miller's
daughter I can understand; but that you should feel embarrassed
before your own first cousin, who has so many times seen you in your
full nakedness~--- that goes altogether too far.

--- Bertel, he replied, I scarcely understand it myself. It is a
curious psychological problem. It is with me as with a fatally ill
person who does not even have the courage to reveal all his dangerous
symptoms to the doctor. In this way a man on many occasions divides
himself into two persons, of whom the one seeks to lead the other
astray, while a third, who is at bottom the same as the other two,
marvels greatly at this confusion. In short, thought becomes dramatic
and plays in silence the most intricate intrigues with itself and for
itself; but the spectator always becomes actor anew. I shall devote a
long chapter to this in my work.

\end{document}
